\documentclass[10pt]{article}
\usepackage[T1]{fontenc}
\usepackage[utf8]{inputenc}
\usepackage[french]{babel}
\usepackage{amsmath,amssymb,mathrsfs,mathpazo}
\usepackage{geometry}
\geometry{a4paper,margin=2cm}
\usepackage{graphicx,array,tabularx,tabularray,colortbl,multirow,multicol,booktabs,enumitem,pifont,ulem,fancybox,diagbox,cancel}
\usepackage{pstricks,pstricks-add,pst-plot,pst-tree,pst-node,pst-text,pst-func,pst-math,pst-bspline,pst-3dplot,pst-eucl,pst-solides3d}
\let\pscalculate\undefined
\usepackage{pst-bezier}
\usepackage{tikz,pgfplots,tkz-tab}
\usetikzlibrary{arrows,arrows.meta,patterns,calc,positioning,trees,decorations.markings,shapes.geometric,shapes.misc,shapes.arrows,intersections,decorations.pathmorphing,decorations.pathreplacing,backgrounds,shadows,babel}
\pgfplotsset{compat=1.18}
\usepgfplotslibrary{fillbetween}
\usepackage[np]{numprint}
\usepackage{listings,makecell,pas-tableur,twoopt}
\let\barre\undefined
\usepackage[pstricks]{tablvar}
\definecolor{darkgreen}{RGB}{0,100,0}
\providecommand{\red}{\color{red}}
\providecommand{\blue}{\color{blue}}
\providecommand{\green}{\color{green}}
\providecommand{\black}{\color{black}}
\providecommand{\white}{\color{white}}
\providecommand{\gray}{\color{gray}}
\providecommand{\cyan}{\color{cyan}}
\providecommand{\magenta}{\color{magenta}}
\providecommand{\yellow}{\color{yellow}}
\providecommand{\violet}{\color{violet}}
\providecommand{\orange}{\color{orange}}
\providecommand{\euro}{\texteuro}
\providecommand{\EUR}{\texteuro}
\providecommand{\ieme}{\textsuperscript{e}}
\providecommand{\degres}{\ensuremath{^\circ}}
\providecommand{\degree}{\ensuremath{^\circ}}
\providecommand{\vv}[1]{\overrightarrow{#1}}
\providecommand{\Coord}[1]{#1}
\providecommand{\decosix}{}
\providecommand{\starredbullet}{\ensuremath{\bullet}}
\pagestyle{empty}
\setlength{\parindent}{0pt}
\usepackage{adjustbox}\DeclareUnicodeCharacter{2236}{:}\DeclareUnicodeCharacter{2190}{\ensuremath{\leftarrow}}
\providecommand{\R}{\mathbb{R}}
\providecommand{\N}{\mathbb{N}}
\providecommand{\Z}{\mathbb{Z}}
\providecommand{\Q}{\mathbb{Q}}
\providecommand{\D}{\mathbb{D}}
\providecommand{\C}{\mathbb{C}}
\providecommand{\e}{{\mathrm{e}}}
\providecommand{\vect}[1]{\overrightarrow{#1}}
\providecommand{\vectt}[1]{\overrightarrow{\mathrm{#1}}}
\providecommand{\vvt}[1]{\overrightarrow{\mathrm{#1}}}
\providecommand{\np}[2][]{{#2~#1}}
\providecommand{\vv}[1]{{\overrightarrow{#1}}}
\renewcommand{\d}{{\mathrm{d}}}
\providecommand{\degres}{{^\circ}}
\providecommand{\degre}{{^\circ}}
\providecommand{\blue}{\color{blue}}
\providecommand{\red}{\color{red}}
\providecommand{\enskip}{\quad}
\providecommand{\llbracket}{\lbrack\!\lbrack}
\providecommand{\rrbracket}{\rbrack\!\rbrack}
\providecommand{\euro}{{\text{€}}}
\providecommand{\no}{{\text{n°}}}
\providecommand{\interval}[2]{{[#1;#2]}}
\providecommand{\intervFO}[2]{{[#1;#2[}}
\providecommand{\dotfill}{\ldots}

\providecommand{\ds}{\displaystyle}
\providecommand{\textding}[1]{#1}
\providecommand{\Oij}{\(\mathrm{O};\vec{\imath},\vec{\jmath}\)}
\providecommand{\Oijk}{\(\mathrm{O};\vec{\imath},\vec{\jmath},\vec{k}\)}
\providecommand{\Ouv}{\(\mathrm{O};\vec{u},\vec{v}\)}
\providecommand{\intervFF}[2]{\ensuremath{[#1;#2]}}
\renewcommand{\euro}{\ifmmode\text{\texteuro}\else\texteuro\fi}
\DeclareUnicodeCharacter{00B0}{\ifmmode^{\circ}\else\textdegree\fi}
\DeclareUnicodeCharacter{20AC}{\ifmmode\text{\texteuro}\else\texteuro\fi}
% Énoncés : Éducation nationale. Transcription et corrigés : APMEP, Denis Vergès et contributeurs. Découpage et classement : CoopMaths, Sébastien Lozano, Éric Elter et contributeurs. Adaptation de la présentation : Maths Explorer.
% CC BY-SA : https://coopmaths.fr/www/about/
% https://coopmaths.fr/alea/static/bac/2025/tex/bac_2025_06_sujet2_metropole_3_cor.tex
\begin{document}
\noindent\textit{APMEP / CoopMaths. CC BY-SA. Présentation adaptée par Maths Explorer.}\par\medskip
\medskip

\begin{enumerate}
\item La suite $(u_n)$ est définie pour tout entier naturel $n$ par
\[u_n = \dfrac{1 + 5^n}{2 +3^n}.\]

\textbf{Affirmation 1 :} La suite $(u_n)$ converge vers $\dfrac 53$. FAUSSE

En écrivant $u_n = \dfrac{1 + 5^n}{2 +3^n} = \dfrac{\frac{1}{5^n} + 1}{\frac{2}{5^n} + \left(\frac 35 \right)^n}$ et en remarquant que :

$\bullet $$\displaystyle\lim_{n \to + \infty} \frac{1}{5^n} = 0$  ;

$\bullet $$\displaystyle\lim_{n \to + \infty} \frac{2}{5^n} = 0$  ;

$\bullet $$\displaystyle\lim_{n \to + \infty} \left(\dfrac 35 \right)^n = 0$ car $0 < \dfrac 35 < 1$ ;

Donc par somme de limites, on a $\displaystyle\lim_{n \to + \infty}\frac{1}{5^n} + 1 = 1$ et $\displaystyle\lim_{n \to + \infty}\frac{2}{5^n} + \left(\frac 35 \right)^n = 0$, donc $\displaystyle\lim_{n \to + \infty} u_n = + \infty$.
\item On considère la suite $(w_n)$ définie par:

\begin{center}$w_0 = 0$\: et, pour tout entier naturel\: $n, w_{n+1} = 3w_n - 2n + 3$.\end{center}

\textbf{Affirmation 2 :} Pour tout entier naturel $n,\: w_n \geqslant n$.

Raisonnement par récurrence :

\emph{Initialisation} : $w_0 = 0 \geqslant 0$  : l'inégalité est vraie au rang 0 ;

\emph{Hérédité} : Soit $n \in \N$ tel que $w_n \geqslant n$, alors 

$3w_n \geqslant 3n \iff 3w_n - 2n \geqslant 3n - 2n$ ou $3w_n - 2n \geqslant n$ et enfin $3w_n - 2n + 3 \geqslant n + 3$.

Or $n + 3 = n + 1 + 2$, donc $n + 3 > n + 1$ et $3w_n - 2n + 3 = w_{n+1}$.

On a donc $w_{n+1}\geqslant n + 3 > n + 1$, soit par transitivité : $w_{n+1}\geqslant n + 1$ : l'inégalité est vraie au rang $n + 1$.

\textbf{Conclusion :} l'inégalité est vraie au rang 0 et si elle est vraie à un rang $n$, elle l'est aussi au rang $n + 1$ ; d'après le principe de récurrence pour tout entier naturel $n,\: w_n \geqslant n$ : l'affirmation 2 est vraie.
\item L'affirmation est fausse car si $f$ était convexe sur $]0~;~+ \infty[$, elle le serait au point d'abscisse 8 et sa courbe représentative serait au dessus de sa tangente au  point A. C'est le contraire d'après le graphique :. Affirmation 3 : fausse.
\item Soit $f$ la fonction définie sur $]0~;~+ \infty[$ par $f(x) = \ln (x) - x + 1$.

Comme somme de fonctions dérivables sur $]0~;~+ \infty[$, $f$ est dérivable sur cet =intervalle et $f'(x) = \dfrac 1x - 1 = \dfrac{1 - x}{x}$.

On a de façon immédiate :

$\bullet~$$f'(x) = 0 \iff x = 1$ ;

$\bullet~$$f'(x) > 0 \iff x < 1$, donc sur ]0~;~1[, \: $f$ est croissante de moins l'infini à $f(1) = \ln 1 - 1 + 1 = 0$.

$\bullet~f'(x) < 0 \iff x > 1$, donc $f$ est décroissante sur $]1~;~+ \infty[$ de 0 à moins l'infini.

Donc cette fonction est négative puisque son maximum est 0 en $x = 1$ : l'affirmation 4 est vraie.
\end{enumerate}

\medskip

\end{document}
